A Theory of Probable Inference — Examples


STATUS

Examples 1, 2 and 3 are built. The rest is the survey pass: every place in the
paper where I think an interactive example earns its keep, numbered in the order
it appears in the text, with the passage it would hang off and what it would do.

BUILD NOTE

The page is now assembled from parts, as the Probability of Induction page is:
run ./build.sh in this folder to write index.html from src/. ./build.sh --check
says whether index.html is behind src/.

  src/01_head.html    <head>, all CSS, opens .article-container
  src/02_article.html Peirce's text, triggers, example containers, footnotes
  src/03_lib.js       maths, the canvas renderer, DOM helpers (from PoI, as-is)
  src/04_scaffold.js  triggers, live numbers, margin numbers (from PoI, as-is)
  src/05_ex123.js     examples 1-3
  src/99_tail.html    closes </body></html>

03 and 04 are the Probability of Induction's own files, unedited, so anything
written for either paper runs in the other. One part file per group of examples
from here on, so an example can be edited without touching the text.

Numbers are assigned in reading order, not build order. (The Errors of
Observation page numbered in build order and left gaps; here the paper is
already whole, so reading order costs nothing and the margin numerals will run
down the page in sequence.) If you would rather number as they are built, say
so — it is one line per example.

The page itself has been reformatted to match the Probability of Induction
edition: same palette, same type, same margin gutter, and the whole example
scaffolding from that page (.example-trigger, .example-container,
.example-header, .ex-num, .live/.k1–k4, .hl-text, .hl-band, .mode-tabs, the row
/ col grid) is already in the stylesheet here. So an example written for that
page drops into this one with no restyling, and the margin number appears beside
its trigger automatically from the container id.

The page-number and paragraph-number marginalia are gone.


GENERAL, CARRIED OVER FROM THE OTHER TWO PAPERS

- Wherever a number or a phrase of Peirce's is being driven, drive it in his own
  sentence: wrap it in .live and colour-code it to the control that moves it.
  Sliders in preference to number entry wherever the quantity is ordered.

- Continuity matters more here than anywhere, because this paper is Peirce
  rebuilding the same three arguments — deduction, induction, hypothesis — from
  five different angles. Where an example restates something the reader already
  met in the Probability of Induction, reuse that example's graphic and say so,
  rather than inventing a second picture of the same thing. Several of the
  examples below are explicitly "PoI's ex-N, re-pointed."

- This paper is more schematic than the other two: its content is very often a
  three-line argument form. So the default graphic here is the argument itself,
  set out, with the terms live and colour-coded, and a population picture beside
  it — not a chart hanging off to one side.

- Peirce's Forms I–V and the "(bis)" variants are already set off with hairlines
  in the page. Those blocks are the natural anchors: an example that belongs to
  a Form should open under that Form's block, not three paragraphs later.

- Let the text do the explaining. Where an example needs a word of its own, say
  it plainly.

- [square brackets] as before: something for the reader to set.


============================================================================
SECTION I — SIMPLE PROBABLE DEDUCTION
============================================================================

## 1

Text: "THE following is an example of the simplest kind of probable inference:—
About two per cent of persons wounded in the liver recover; This man has been
wounded in the liver; Therefore, there are two chances out of a hundred that he
will recover." — through the Enoch syllogism and Forms I and II.

Suggestion. The whole point of the opening is that Forms I and II are the same
argument with one knob turned, so build one argument and turn the knob. A single
slider for [rho], from 0 to 1, sitting under Form II. At the left of the track a
hundred figures, [rho] of them marked as recoveries; the conclusion sentence
under it reads back the chance. As [rho] runs to 1 the block goes solid, the
premise reads "Every M is a P", and Form II's three lines visibly become Form
I's — same layout, the words swapping in place. The liver figure (2 per cent)
and the Enoch figure (100 per cent) are two marked stops on the same slider, so
the reader can flip between Peirce's two examples without leaving the control.

The thing to make unmissable: nothing about the *form* changes as [rho] moves.
Only the strength does. That is the claim of the section and it is made in one
gesture.

[great]

--- BUILT: EXAMPLE 1 --------------------------------------------------------

Margin number: 1.  Trigger: "THE following is an example of the simplest kind of
probable inference:—".  Container: under Form II.  Source: src/05_ex123.js.

What it does. One slider, rho, from 0 to 1, with Peirce's two figures marked on
it as buttons (2% and 100%) plus an even chance and 98%. Left of the column, a
hundred squares with rho of them scattered as P's; right of it, the argument
twice over — the general form and the liver wound — three lines each, aligned,
with the live figures in the slider's blue.

At 100 the major premise turns into "Every M is a P", the conclusion loses its
"with probability", the block's rule goes green, and the name above it reads
Form I, Singular syllogism in Barbara. At 0 it does the same negatively. Between
them it is Form II throughout. The two live spans in Peirce's own scheme at the
head of the section move with it, so his sentence reads "About fifty per cent of
persons wounded in the liver recover... there are fifty chances out of a
hundred" while you drag.

Underneath: draw this man / draw fifty, which tallies how often the argument
came out true. That is the "strength" the slider is really moving — the form is
the same whichever man you draw.

Open questions for you:
- The hundred squares are a fixed scatter, not a filling bar, so 2% reads as two
  people here and there. Say if you would rather see it fill.
- The tally resets whenever rho moves. Keeping it would let you compare rates
  across settings, but it would also be comparing two different arguments.


## 2

Text: "It is to be observed that the ratio ρ need not be exactly specified. We
may reason from the premise that not more than two per cent of persons wounded
in the liver recover, or from 'not less than a certain proportion of the M's are
P's,' or from 'no very large nor very small proportion, etc.' In short, ρ is
subject to every kind of indeterminacy; it simply excludes some ratios and
admits the possibility of the rest."

Suggestion. A 0–1 bar for the value of ρ, on which the reader paints the
admitted region — a two-handled range for "not more than 2 per cent" or "no very
large nor very small proportion", and a scatter of admitted/excluded bands for
the messier cases. Buttons for Peirce's three phrasings, plus a free one the
reader draws. Under it, the conclusion, which is now itself a region and not a
number: "it follows, with probability [between .00 and .02], that S is a P."

This is the first appearance of a thing the paper keeps needing — a probability
that is a set of values rather than one value — so it is worth spending a
graphic on it early, and reusing that bar later (16, 22) rather than drawing a
new one.

[great]

--- BUILT: EXAMPLE 2 --------------------------------------------------------

Margin number: 2.  Trigger: the whole "It is to be observed that the ratio rho
need not be exactly specified" passage.  Source: src/05_ex123.js.

What it does. A 0-1 bar with a painted band, driven by two handles, from and to.
Five buttons: exactly two per cent (the band collapses to a line), not more than
two per cent, not less than…, no very large nor very small, and anything but
middling. A checkbox turns the band into what the premise EXCLUDES, which is how
the two-outer-bands cases are reached.

Peirce's three phrasings are marked up in his own sentence and light in three
colours as you pick the button that matches them — the same .hl-text treatment
the other paper uses for antecedent / consequent / consequence.

Below the bar, Form II with the ratio left indeterminate: both the premise and
the conclusion read "between 0.15 and 0.85", so the conclusion is visibly a
region of the same bar rather than a point on it. The note counts how many of
the hundred ratios the premise admits and how many it excludes, and says so in
his words. Push the band to the full width and the note says the premise now
excludes nothing, which is to say there is no conclusion.

Open questions for you:
- Two sliders rather than one two-handled control. A real dual-range needs hand
  -rolling; say if it is worth it.
- The bar is the one I proposed reusing at 15 and 22. If you want it to look
  different there, better to say now.


## 3

Text: "Thus, there being four aces in a picquet pack of thirty-two cards, the
chance is one eighth that a given card not looked at is an ace; but this is only
on the supposition that the card has been drawn at random from the whole pack.
If, for instance, it had been drawn from the cards discarded by the players at
piquet or euchre, the probability would be quite different." — and, four
paragraphs on, "But after we have looked at the card, we can no longer reason in
that way."

Suggestion. Reuse the card grid from the Probability of Induction (examples 1–4
there — createGrid / evaluateConsequent already handle the piquet pack). Deck
choice: whole piquet pack, or the discards. Draw a card face down; the grid
beside it shades the cells still live and reports the chance. Then a "look at
the card" button, which turns it face up and collapses the grid to a single
cell — the chance is now 1 or 0, and the argument the reader was making a moment
ago is not available any more.

Two conditions of the same maxim in one control, which is how Peirce presents
them: the drawing has to be random, and the conclusion has to be drawn in
advance of any other knowledge. Worth colour-coding the two premises in his
sentence to the two halves of the control.

[great, lets also have a setting for making the cards discarded from a game, so they would appear as hands]

--- BUILT: EXAMPLE 3 --------------------------------------------------------

Margin number: 3.  Triggers: two — "Thus, there being four aces in a picquet
pack…", where the example sits, and "But after we have looked at the card, we
can no longer reason in that way" four paragraphs down, which opens the same
one and scrolls to it.  Source: src/05_ex123.js.

What it does. Three sources for one card, as buttons:

 - the whole piquet pack: 32 cards drawn as cards, the four aces outlined.
   Chance 4/32, and the live span in Peirce's sentence reads "one eighth".
 - the cards discarded at piquet: two hands of twelve are dealt and each player
   throws out five. Bright cards are the discards, dimmed ones are kept; the
   pile of ten is drawn out below them. This is your "so they would appear as
   hands".
 - the cards discarded for euchre: the full 52 with the euchre pack dimmed, so
   the pile to draw from is the twos to eights — and every ace is in the other
   half.

The piquet case got a knob rather than a fixed strategy: how far the players
choose what to throw out. At 1 they always throw the five lowest and no ace ever
reaches the pile; at 0 they discard at random and the pile is as good as the
pack, an ace or so in ten. Dragging it walks the chance from one eighth down to
nothing, which is the maxim as a dial rather than a switch — the will
intermeddling by degree.

Then the second condition. Draw a card (it comes face down, and the argument
reads normally), then LOOK AT IT: the card flips, the conclusion line goes red
and says the premises have nothing left to do, and the live figure in Peirce's
sentence four paragraphs down changes with it. Deal again and the eighth
returns.

Open questions for you:
- The euchre reading is that "the cards discarded by the players at euchre" are
  the low cards thrown out of a full pack to make the 24-card euchre pack. The
  other reading is the dealer's discard after turning up the trump, which is one
  card and hard to draw from. Tell me if you meant the other one.
- The care dial defaults to 1, so the first thing you see is the strongest
  version of his point. It could default to something middling instead.


## 4

Text: "The celebrated law of Fechner is, that as the force acting upon an organ
of sense increases in geometrical progression, the intensity of the sensation
increases in arithmetical progression. In this case the odds ... take the place
of the exciting cause, while the sensation itself is the feeling of confidence.
When two arguments tend to the same conclusion, our confidence in the latter is
equal to the sum of what the two arguments separately would produce; the odds
are the product of the odds in favor of the two arguments separately."

Suggestion. This is the Probability of Induction's example 26 (the logarithm of
the chance) arriving from the other direction, and the graphic should be
recognisably the same one: the odds on one scale, the feeling of belief on the
log scale beside it, null at odds 1/1. Add what this passage has that the other
did not — two arguments, each with its own odds, whose chances multiply while
their beliefs add. Two sliders, two bars, and the sum/product shown on both
scales at once so the reader sees the same fact twice.

If the reader has come from the other paper, say in a line that this is the same
thermometer.

[great]

============================================================================
SECTION II — COMPLEX AND STATISTICAL DEDUCTION
============================================================================

## 5

Text: Form III (Complex Probable Deduction), the factorial formula for q, and
"As an example, let us assume the proportion r = 2/3 and the number of M's in a
set n = 15. Then the values of the probability q for different numbers, m, of
P's, are fractions having for their common denominator 14,348,907..."

Suggestion. The table of numerators is on the page already; make it the example.
Sliders for [r] and [n], both live in Peirce's sentence; the table redraws, and
the common denominator with it. Beside it the same numbers as a bar chart, with
the bar at the maximum marked and the value (n+1)r drawn as a rule across the
chart so the reader can watch the claim in the next paragraph — that q peaks at
the m next below (n+1)r, and is very small elsewhere — hold as they drag.

Footnote 1's case (when (n+1)r is a whole number, two bars tie) should be
reachable: at r = 1/2, n = 15 it happens, and the two equal bars are worth
seeing rather than reading.


[that sounds good, but we should visualize the n sets of Ms.  and the proportion (with colour) we shoudl show the sample and the drawing, so we see the Ss from the Ms. please double check the wording of this example in the html against the text]. 

## 6

Text: "When this is the case, there is a more convenient way of obtaining (not
exactly, but quite near enough for all practical purposes) either a single value
of q or the sum of successive values..." — the t1 / t2 formulae and the table
of Θt.

Suggestion. Same chart as 5, now with a shaded band from [m1] to [m2]. Two
readouts side by side: the exact sum of the q's in the band, and Peirce's
½Θt2 − ½Θt1. The Θ table on the page lights the two rows being read, the way
the probable-error table does in the other paper's example 17.

The point to land is "quite near enough": show the difference between the two
numbers, and let the reader push n up until it stops mattering. His rough rule —
Θt ≈ t below 0.7, Θt ≈ 1 above 1.4 — can be a third readout that goes green when
it is safe to use.


## 7

Text: Form IV (Statistical Deduction) and Form V (Induction), and the paragraph
between them: "If, then, this principle justifies our inferring the value of the
second proportion from the known value of the first, it equally justifies our
inferring the value of the first from that of the second, if the first is
unknown but the second has been observed."

Suggestion. The hinge of the paper, and it deserves the biggest example in it.
One population, one sample, two arrows. Mode tabs: DEDUCTION runs the arrow
downward — set [r] in the population, draw [n], watch the sample proportions
scatter about r. INDUCTION runs it upward — the population is hidden, the reader
draws a sample and the inference points back at a region of the population bar.
Same graphic, same numbers, arrow reversed; nothing else changes, which is
exactly Peirce's argument for why their validity is the same.

Then the second half of the passage, which is the part people miss: keep drawing
in each mode. In deduction the prediction is *vindicated* — the running
proportion settles onto the r you set. In induction it is *corrected* — the
inferred region moves off the first guess and closes on the truth. Two little
running-trace plots, one per mode, both labelled in his words.

Peirce's own instances belong here as preset buttons: the bag of coffee (nine
tenths of the beans perfect) and the 1870 census births (478,774 to 463,320
white, 75,985 to 76,637 colored). The census one is the same data as example 16
in the Probability of Induction — link to it rather than re-deriving the
probable errors here.


============================================================================
SECTION IV — DEDUCTION AND HYPOTHESIS "IN DEPTH"
============================================================================

## 8

Text: "We often speak of one thing being very much like another, and thus apply
a vague quantity to resemblance... We may then measure resemblance by a scale of
numbers from zero up to unity. To say that S has a 1-likeness to a P will mean
that it has every character of a P" — through Form II (bis) and Form IV (bis),
and the French/Italian example.

Suggestion. The whole trick of this section is that "in depth" is the same
picture with characters on the axis instead of individuals, so build the mirror
of 7 and let the reader see the axis swap. A strip of marks P′, P″, P‴… instead
of a row of objects; a slider for [r]-likeness; the marks S shares light up.
At r = 1 every mark is shared and the form reduces to Barbara — same collapse as
example 1, and worth animating the same way for the rhyme.

For Form IV (bis), Peirce's boy travelling through France and Italy: the "really
numerous though relatively few respects in which he will be able to compare the
two people" is a random sample of the marks, and the shared proportion in the
sample tracks r. That is a two-panel version of the same strip — the full run of
marks and the few he actually gets to compare.

A toggle marked "count them / weigh them" belongs somewhere here or in 9, for
the objection Peirce raises himself on the next page: characters cannot be
counted like individuals — antimony bluish-gray, bismuth rose-gray, and gold,
silver, copper, tin further off. Give the marks weights and the proportion
becomes a weighted one; his own complaint about his own form, made visible.


## 9

Text: Form V (bis), Hypothesis, and "Thus, we know, that the ancient
Mound-builders of North America present, in all those respects in which we have
been able to make the comparison, a limited degree of resemblance with the
Pueblo Indians. The inference is, then, that in all respects there is about the
same degree of resemblance between these races."

Suggestion. Invert 8, exactly as 7 inverts its own first mode, and reuse its
graphic: the marks compared so far are the sample, the r-likeness inferred is
the conclusion, the full run of characters is hidden. Draw more comparisons and
watch the inferred likeness get corrected rather than vindicated.

The three inferences now all exist as the same picture at two axes and two
directions. Worth a small four-cell diagram, once, showing where the reader has
been: breadth/depth × deduction/ampliation, with Peirce's Form numbers in the
cells and each cell clicking back to its own example.


============================================================================
SECTION V — DEDUCTION, INDUCTION, HYPOTHESIS DISTINGUISHED
============================================================================

## 10

Text: "If I wished to order a font of type expressly for the printing of this
book, knowing, as I do, that in all English writing the letter e occurs oftener
than any other letter, I should want more e's in my font than other letters...
This is a statistical deduction. But then the words used in logical writings are
rather peculiar... I might, then, count the number of occurrences of the
different letters upon a dozen or so pages of the manuscript... That would be
inductive inference. If now I were to order the font, and if, after some days, I
were to receive a box containing a large number of little paper parcels of very
different sizes, I should naturally infer that this was the font of types I had
ordered; and this would be hypothetic inference."

Suggestion. The best passage in the paper for an example, because all three
inferences run on one body of evidence, and the body of evidence can be this
page. Reuse the letter-reading machinery from the Probability of Induction's
vowel/consonant example (the .vc spans) over the text of *this* article.

Three tabs on one letter-frequency chart:
 - DEDUCTION: the standard English frequencies, laid over the chart as a rule.
   Order the font from them.
 - INDUCTION: read [a dozen or so] pages of the actual article, letter by
   letter, watching the observed bars climb and depart from the English rule —
   Peirce is right, logical writing is peculiar; the single letters used as
   terms show up as a visible bulge.
 - HYPOTHESIS: the box of parcels. Show only the parcel sizes — the bars with
   their labels stripped off — and let the reader match that shape back to the
   font. The point is that the characters are observed and the case is inferred.

One dataset, three directions of inference, and the reader has watched the
difference instead of being told it.


## 11

Text: "Again, if a dispatch in cipher is captured, and it is found to be written
with twenty-six characters, one of which occurs much more frequently than any of
the others, we are at once led to suppose that each character represents a
letter, and that the one occurring so frequently stands for e. This is also
hypothetic inference."

Suggestion. A short ciphered dispatch on screen with its frequency bars beside
it. Click the tallest bar and assign it e; the substitution propagates through
the text and partial words appear. Two or three more clicks and it reads.

Keep it small — it is a rider on 10, not a cryptogram puzzle. What it has to
show is the shape of the inference: the rule (English has these frequencies)
plus the result (this text has that profile) gives the case (it is English in
substitution). Set the three lines out beside the bars in the same Rule / Case /
Result type used further down the section.


## 12

Text: "Now in Barbara we have a Rule, a Case under the Rule, and the inference of
the Result of that rule in that case. For example:— Rule. All men are mortal;
Case. Enoch was a man. Result. Enoch was mortal." — and "Deduction proceeds from
Rule and Case to Result... Induction proceeds from Case and Result to Rule...
Hypothesis proceeds from Rule and Result to Case."

Suggestion. Three propositions on a triangle; the reader covers one and the
other two infer it. Cover the Result — deduction. Cover the Rule — induction.
Cover the Case — hypothesis. Each cover changes the colour of the arrows and
names the inference in Peirce's words underneath, with the psychological gloss
he attaches (volition, the formation of a habit, secondary sensation) as a
second line the reader can turn on.

This is the paper's central mnemonic and it is three clicks. It should also be
the thing that examples 7, 9, 10 and 11 all link back to, so the reader can see
which corner each of them was covering.


============================================================================
SECTION VI — THE RULE, AND APAGOGIC INVERSION
============================================================================

## 13

Text: "In the case of hypothesis, this syllogism is called the explanation...
The explanation is,— Simple English ciphers have certain peculiarities; This is
a simple English cipher: Hence, this necessarily has these peculiarities." —
and the barrel of apples, "not 'Why is this?' but 'How is this?'"

Suggestion. Take the reader's cipher from 11 and the apples together, side by
side, and run the same three-line explanatory syllogism under each: the Why of
hypothesis and the How of induction, identical in structure, different only in
which premise is being recovered. Sliders on the barrel — [proportion decayed]
and [size of the handful] — with the sample drawn and the explanation written
out from it in Peirce's words.

Then the test the rule actually imposes: a switch that makes the explanatory
syllogism a *bad* statistical deduction (draw the handful off the top of the
barrel, or predesignate nothing), and the inference above it greys out. The rule
of section VI is one rule, and this is the place to make it operable rather than
stated.


## 14

Text: "The statistical syllogism of Form IV. is invertible, because it proceeds
upon the principle of an approximate equality between the ratio of P's in the
whole class and the ratio in a well-drawn sample, and because equality is a
convertible relation. But ordinary syllogism is based upon the property of the
relation of containing and contained, and that is not a convertible relation."
— through the three figures and the three modi, now set out on the page.

Suggestion. Two little relation diagrams, side by side, each with a "reverse it"
button. Equality reverses and stays true. Containment reverses and breaks —
show the counter-case as two circles, one inside the other, and the false
converse written under them.

Then the schemes: the three figures and the three moods are on the page as
columns, and the transformation between them is mechanical — take the
conclusion and one premise, swap them, negate both. Make that an animation the
reader steps through, one figure at a time, with the negated lines flipping in
place. This is also the machinery the missing pages would have applied to
Form IV, so the example partly stands in for what the scan lost; the editorial
note about the gap sits right below it and can say so.


============================================================================
SECTION VII — DRAWING AT RANDOM
============================================================================

## 15

Text: "Suppose that, instead of using such a precept of selection that any one M
would in the long run be chosen as often as any other, we used a precept which
would give a preference to a certain half of the M's, so that they would be
drawn twice as often as the rest... we could deduce by arithmetic the further
conclusion that, counting the M's for one each, the proportion of P's among them
must (ρ being over 2/3) lie between 3/4 ρ + 1/4 and 3/4 ρ − 1/2."

Suggestion. A population split into a preferred half and the rest, with a slider
for [how much oftener] the preferred half gets drawn — Peirce's case is twice,
but the general knob is more instructive. Sample from it, get ρ, and show two
readings side by side: the weighted conclusion (which is exactly right), and the
unweighted one, which is now the interval 3/4 ρ + 1/4 to 3/4 ρ − 1/2 painted on
the ratio bar from example 2.

Then run ρ up towards 1 and down towards 0 and watch that interval close, which
is his consoling remark: "when ρ approximates towards 1 or 0, the effect of the
imperfect sampling will be but slight." The bad sampling does not destroy the
inference, it blunts it, and the reader can see how much.


## 16

Text: "Nor must we lose sight of the constant tendency of the inductive process
to correct itself. This is of its essence. This is the marvel of it... a
different selection, made by a different method, will be likely to vary from the
normal in a different way, and if the ratios derived from such different
selections are nearly equal, they may be presumed to be near the truth." — and
the warning after footnote 6, that "the magnitude of its probable error will
depend very much more on the worst than on the best inductions used."

Suggestion. Three or four sampling precepts running down the page, each biased
in its own direction, each producing its own inferred interval. The reader adds
instances and watches all of the intervals move; where they overlap, a combined
reading. Then a control that makes one precept much worse than the others, and
the combined reading degrades with it — not to the average of the four, but
towards the worst. Two claims, one picture, and the second is the one that gets
forgotten.


## 17

Text: "Within the period of time M, a certain event P occurs; S is a period of
time taken at random from M, and more than half as long: Hence, probably the
event P will occur within the time S." — and its inversion, and "This is a sort
of logic which is often applied by physicists in what is called extrapolation of
an empirical law."

Suggestion. A timeline for M with the event P somewhere on it, and a window S
the reader drags and stretches. Above half the length of M, the window catches P
more often than not — run it a thousand times and show the hit rate against the
window length, with the half-length mark on the axis. Then the inversion: the
window is the history available, P did not happen in it, conclude it does not
happen in M.

Worth ending on his own verdict rather than on the success rate: "as compared
with a typical induction, it is obviously an excessively weak kind of
inference." Show the hit rate at just over half and let it be unimpressive.


## 18

Text: "Conceive a cardboard disk revolving in its own plane about its centre,
and pretty accurately balanced, so that when put into rotation it shall be about
as likely to come to rest in any one position as in any other; and let a fixed
pointer indicate a position on the disk: the number of points on the
circumference is infinite, and on rotating the disk repeatedly the pointer
enables us to make a selection from this infinite number."

Suggestion. Spin the disk. Each stop drops a mark on the circumference and adds
a value to a list, and the list never repeats. Then the second half of the
passage, which is the surprising one: a finite urn sampled with replacement is
the same thing — what is really being sampled is "not the finite collection of
things, but the unlimited number of possible drawings." Put the urn beside the
disk and run both; two counters, both unbounded.

Small and cheap, and it sets up the footnote-8 identity in 19.


## 19

Text: "But in what is known as 'reasoning from analogy,' the class sampled is
small, and no instance is taken twice. For example: we know that of the major
planets the Earth, Mars, Jupiter, and Saturn revolve on their axes, and we
conclude that the remaining four, Mercury, Venus, Uranus, and Neptune, probably
do the like." — through the three-premise scheme, the paired syllogisms, and
footnote 8.

Suggestion. Eight planets, four observed. The three premises set out live: the
sample of objects, the sample of characters, the conclusion. Then split the
argument into the hypothesis and the induction that compose it — the page
already prints them side by side — and light each half as the reader hovers it,
with the deduction falling out below.

Footnote 8 is the interesting claim and it can be checked on screen: drawing
without replacement from a small class gives the same chances as drawing from an
infinite class that the small class was itself drawn from. Two urns, the two
expressions, and a run counter that converges to the same number. That is the
justification for treating four planets as a sample of anything at all, and it
is currently buried in a footnote.


============================================================================
SECTION VIII — PREDESIGNATION
============================================================================

## 20

Text: "Take any human being, at random,—say Queen Elizabeth. Now a little more
than half of all the human beings who have ever existed have been males; but it
does not follow that it is a little more likely than not that Queen Elizabeth
was a male, since we know she was a woman."

Suggestion. Short and sharp. Draw a person at random from a population bar that
is a little over half male; the inference reads "probably male". Then a button:
"look at who you drew." The precept has changed because the case presented to
the mind has changed, and the inference is not available. Same shape as the
face-up card in example 3, deliberately — this is Peirce making the same point
about knowledge in the second premise, and the rhyme is worth the reuse.


## 21

Text: "I take from a biographical dictionary the first five names of poets, with
their ages at death... These five ages have the following characters in
common:— 1. The difference of the two digits composing the number, divided by
three, leaves a remainder of one..." — and "Yet there is not the smallest reason
to believe that the next poet's age would possess these characters."

Suggestion. The best joke in the paper and it should be playable. A button that
draws [five] ages at random, and a machine that then hunts for characters they
all share — digit arithmetic, prime factors, residues mod 3, the same family of
tricks Peirce used — and prints the two or three it finds. Draw again: new ages,
new characters, always some. A counter for how many draws in a row the machine
has succeeded (it will not fail).

Then the other half, which is what makes it an argument and not a trick:
predesignate. The reader picks a character *first*, from the same list, and then
draws. Now the success rate is what it should be. Two counters side by side,
running, is the whole doctrine of predesignation.

His demonstration that this must always be possible — a curve can be drawn
through any given series of points while missing any other given series — could
be a small second panel: points on a plane, a curve through the chosen ones,
avoiding the rest, redrawn as the reader adds points.


## 22

Text: Playfair on the specific gravities of the allotropic forms of carbon,
through "It thus appears that there is no more frequent agreement with
Playfair's proposed law than what is due to chance."

Suggestion. The tables are already on the page; make them the example. Stage
one: Playfair's own progression — carbon, then silicon (needing a modified
formula), then boron (needing another), then phosphorus, sulphur, selenium,
tellurium, bromine, iodine. Each element arrives as a row with a tick or a
cross, and a running count of predesignate cases: 8, of which 4 agree.

Stage two is Peirce's test and it is a beautiful piece of statistics for 1883:
pair each element's specific gravity with the atomic weight of a *different*
element, and count the agreements again. The page prints his pairing; let the
reader shuffle the pairing themselves and watch the hit rate stay where it was.
A histogram of hit rates over many random re-pairings, with Playfair's own
result marked on it, is a permutation test in everything but name, and the mark
sits right in the middle of the distribution.

Flag in a note that this is what the modern reader would call a permutation
test — but plainly, one sentence, no fanfare.


## 23

Text: "Suppose that, before sampling a class of objects, we have predesignated
not a single character but n characters, for which we propose to examine the
samples. This is equivalent to making n different inductions from the same
instances. The probable error in this case is that error whose probability for a
simple induction is only (1/2)^n, and the theory of probabilities shows that it
increases but slowly with n; in fact, for n = 1000 it is only about five times
as great as for n = 1..."

Suggestion. A slider for [n], the number of characters predesignated, and two
readouts: the probable error, and the number of instances needed to buy it back.
The shape of the curve is the message — it climbs steeply for the first few and
then almost flattens, so 25× the instances covers n = 1000 and 100× covers ten
billion. Mark on the same axis the sizes Peirce says actually occur: "of
characters at all striking, or of objects at all familiar, the number will
seldom reach 1,000."

The conclusion the reader should reach by dragging, before reading it: failing
to predesignate is a real fault, and a surprisingly cheap one when the instances
are many. And it is the same slider that explains why "a hypothesis ought to be
simple" — simple here meaning familiar, which is to say drawn from a small n.


============================================================================
SECTION IX — ANTECEDENT KNOWLEDGE
============================================================================

## 24

Text: "There are thus four different kinds of uniformity and non-uniformity which
may influence our ampliative inferences:— 1. The members of a class may present
a greater or less general resemblance as regards a certain line of characters..."
— with the gold, the three-toed genus, the colour that varies in almost every
genus, and the Icelanders.

Suggestion. Four knobs, one per kind, over the same population-and-sample
picture used in 7, so the reader can feel each one separately widen or narrow
the conclusion drawn from a fixed sample. Peirce's own instances as presets:
gold's chemical character (kind 1, very high — two or three specimens settle
it), three toes on each foot in a genus (kind 2, generic character), colour in a
genus (kind 2, very low), the Icelanders (kind 1 again, but for ideas), the
purple of cassius (kind 4, a characteristic test), the man who carries every
opinion to its logical end (kind 4).

The end state to reach: turn a knob far enough and the inference stops being
ampliative at all — "the inference is mere deduction, that is, the application of
a general rule already established." The example should be able to go all the
way there and say so.


============================================================================
SECTION X — AGAINST INVERSE PROBABILITY
============================================================================

## 25

Text: "Suppose an ancient denizen of the Mediterranean coast, who had never
heard of the tides, had wandered to the shore of the Atlantic Ocean, and there,
on a certain number m of successive days had witnessed the rise of the sea.
Then, says Quetelet, he would have been entitled to conclude that there was a
probability equal to (m+1)/(m+2)..." — and "Putting m = 0, it is seen that this
view assumes that the probability of a totally unknown event is 1/2."

Suggestion. This is the Probability of Induction's example 29 almost word for
word, and it should be visibly the same example: same shore, same slider for
[m], same (m+1)/(m+2) reading out in the sentence. Reuse it wholesale, with the
m = 0 stop marked, since the whole force here is what happens at zero.

What this paper adds, and the earlier one does not have: "In point of fact, we
know that although theories are not proposed unless they present some decided
plausibility, nothing like one half turn out to be true." Worth a second panel —
a rate the reader can set for how often proposed theories are true, and the
rule of succession's assumed one-half sitting there, wrong, next to it.


## 26

Text: "Suppose, for instance, there were four possible occasions upon which an
event might occur. Then there would be 16 'constitutions of the universe,' or
possible distributions of occurrences and non-occurrences." — both tables, and
"assuming the different frequencies of the first event to be equally probable,
those of this new event are not so."

Suggestion. Both tables are on the page. Make the sixteen strings the object:
a switch between EQUAL FREQUENCIES and EQUAL CONSTITUTIONS, which reweights the
sixteen cells visibly — under the first, the five columns get equal weight and
the six-member column's cells are each worth a sixth of it; under the second,
every cell is worth the same.

Then Peirce's actual argument, which is the derived event: define the new event
as "a Y following a Y or an N following an N", let the reader watch each of the
sixteen strings sort itself into the second table, and read off the induced
weights. Under equal constitutions they stay equal; under equal frequencies they
do not — three occurrences comes out half as large again as two, exactly as he
says. The inconsistency is a consequence the reader derives by flipping one
switch, which is much better than being told it.

Then the sting, and it wants stating in the example: adopt equal constitutions
and "the occurrences or non-occurrences of an event in the past in no way affect
the probability of its occurrence in the future." Let the reader run draws under
that assumption and see the past teach them nothing.


## 27

Text: "Suppose, for example, that I were to enter a great hall where people were
playing rouge et noir at many tables; and suppose that I knew that the red and
black were turned up with equal frequency... But could some devil look at each
card before it was turned, and then influence me mentally to bet upon it or to
refrain therefrom, the observed ratio in the cases upon which I had bet might be
quite different from the observed ratio in those cases upon which I had not
bet."

Suggestion. A hall of tables, red and black turning up evenly. The reader bets
at will and the two ratios — bet-on and not-bet-on — are tracked side by side.
They agree. Then hand the devil a switch: he sees each card and nudges the
choice. The two ratios come apart, and both the deduction and the induction
built on them break at the same moment.

That simultaneity is the point of the passage, and it is why the example is
worth building rather than summarising: Peirce is conceding that his theory
needs a fact supposed too, and then showing that the fact it needs is only the
absence of a devil — the same thing deduction needs. Two panels, one devil, both
break together.


============================================================================
SECTION XI — THE LOGIC OF SCIENTIFIC INVESTIGATION
============================================================================

## 28

Text: "Suppose a being from some remote part of the universe... to be presented
with a United States Census Report... He begins, perhaps, by comparing the ratio
of indebtedness to deaths by consumption in counties whose names begin with the
different letters of the alphabet." — through the January rain-fall and
illiteracy table and the two plates.

Suggestion. The last example, and it should be the reader taking the stranger's
seat. A census-like table of columns, and a pair of dropdowns: compare any
column against any other. Most pairs give a flat line and the readout says, in
Peirce's terms, that the proportions among the M's and the non-M's do not
differ — the induction is unimportant. Let the reader grind through several of
those, including his own joke pairing by initial letter, before finding
anything.

Then the pairing he does find: the seven regions plotted, January rain-fall
against illiteracy, with the two maps already on the page beside it. It is a
real but unsatisfactory correlation and he says so — "it provokes further
inquiry; we desire to replace the M by some different class, so that the two
proportions may be more widely separated." Give the reader a way to split the
regions further and watch the separation improve or not.

The moral is the one the page ends on, and the example should hand off to it
rather than restate it: the stranger cannot guess which pairs to try, and we
can. Guessing right is the faculty, not the arithmetic.


============================================================================
NOT WORTH AN EXAMPLE, AND WHY
============================================================================

- The whole of the discussion of the copula of inclusion (section V, "This
  classification of probable inference is connected with a preference for the
  copula of inclusion over those used by Miss Ladd and by Mr. Mitchell"). It is
  a remark to his co-authors in Studies in Logic. Nothing moves.

- The decapitated frog. It is an analogy, not a mechanism, and a picture of it
  would be either grisly or coy.

- Section X's closing pages on Gratry and the miraculous intervention, and on
  whether the mind is adapted to the universe. Argument, not quantity.

- The gap at pp. 150–151. Tempting to reconstruct Form VI in a "coming soon"
  block, since example 14 gets close to it, but reconstructing missing Peirce
  and presenting it interactively is a different kind of act from illustrating
  him. Left alone unless you want it.


============================================================================
PROGRESS LOG
============================================================================

(One block per example as it is built, appended here, per the convention in the
Errors of Observation notes.)

1  built.  Awaiting your read.
2  built.  Awaiting your read.
3  built.  Awaiting your read; two open questions above, on euchre and on the
   default setting of the care dial.
4-28  not built.
